\documentclass{ximera}

\title{Oefeningen Elektrisch Veld}
\author{Daan Lipkens}

\begin{document}
\begin{abstract}
    Test je inzicht in en rekenvaardigheid over het elektrisch veld.
\end{abstract}
\maketitle



\section{Vraagstukken}

%———————————————————————————————————————
% VRAAG: Problem 12 vertaald
\begin{exercise}\label{vb:elektrisch_veld_2ladingen}
    Twee ladingen, $Q_1 = -2,50 \,\mu\mathrm{C}$ en $Q_2 = 6,00 \,\mu\mathrm{C}$, bevinden zich op een horizontale as, gescheiden door een afstand van $1,00\,\mathrm{ m}$. Bepaal het punt op de as (anders dan het oneindige) waar het netto elektrisch veld precies nul is.
    

        \begin{hint}
            Tussen de twee ladingen in wijzen de velden van beide ladingen in dezelfde richting. Het veld kan daar dus nooit nul worden. Het punt moet aan de \textit{buitenkant} liggen.
        \end{hint}
        \begin{hint}
            Het nulpunt ligt altijd aan de buitenkant, aan de zijde van de lading met de \textit{kleinste absolute waarde}. Dat is in dit geval $Q_1$ ($-2,50\,\mu\mathrm{C}$).
        \end{hint}
        \begin{solution}\nl
            Noem de afstand vanaf de lading $q_1$ tot het nulpunt $d$. De afstand tot $q_2$ is dan $(d + 1,00\,\mathrm{m})$. We stellen de grootte van de velden aan elkaar gelijk:
            $$ E_1 = E_2 $$
            $$ k_e \frac{|Q_1|}{d^2} = k_e \frac{|Q_2|}{(d + 1,00 \,\mathrm{m})^2} $$
            Deelt we $k_e$ en de $\mu C$-factor weg en nemen we vierkantswortels, dan krijgen we:
            $$ \frac{\sqrt{2,50}}{d} = \frac{\sqrt{6,00}}{d + 1,00\,\mathrm{m}} $$
            Na wiskundige omvorming levert dit een afstand op van $d = 1,82\,\mathrm{m}$ aan de linkerkant van de lading $Q_1$.
        \end{solution}
\end{exercise}
%———————————————————————————————————————

% %———————————————————————————————————————
% % VRAAG: Problem 36 vertaald
% \begin{exercise}
%     Een proton wordt op tijdstip $t = 0\text{ s}$ horizontaal (in de positieve x-richting) afgeschoten in een homogeen elektrisch veld met veldvector $\vec{E} = -6,00 \cdot 10^5 \cdot \hat{i}\text{ N/C}$. Het proton legt een afstand van $7,00\text{ cm}$ af vooraleer het volledig tot stilstand komt.
    
%     Gegeven: $m_p = 1,67 \cdot 10^{-27}\text{ kg}$ en $q_p = 1,60 \cdot 10^{-19}\text{ C}$.
    
%     Bereken:
%     \begin{enumerate}
%         \item De \textbf{versnelling} van het proton.
%         \item De \textbf{beginsnelheid} van het proton.
%         \item Het \textbf{tijdsinterval} tot stilstand.
%     \end{enumerate}
    
    
%         \begin{hint}
%             Gebruik de tweede wet van Newton $\vec{F}_e = m\vec{a}$ waarbij $\vec{F}_e = q\vec{E}$.
%         \end{hint}
%         \begin{hint}
%             Omdat het veld constant is, mag je voor de kinematica de formules voor een eenparig veranderlijke rechtlijnige beweging (EVRB) gebruiken. Let erop dat het veld in de \textit{negatieve} zin werkt, dus het proton ondervindt een remmende kracht!
%         \end{hint}
%         \begin{solution}\nl
%             \textbf{1. De versnelling:}
%             $$ a = \frac{qE}{m} = \frac{(1,60 \cdot 10^{-19}\text{ C}) \cdot (-6,00 \cdot 10^5\text{ N/C})}{1,67 \cdot 10^{-27}\text{ kg}} = -5,75 \cdot 10^{13}\text{ m/s}^2 $$
%             De versnelling is negatief en werkt dus remmend.
            
%             \textbf{2. De beginsnelheid:}
%             Met EVRB-formule (zonder tijd): $v_f^2 = v_i^2 + 2a\Delta x$
%             $$ 0 = v_i^2 + 2(-5,75 \cdot 10^{13}\text{ m/s}^2)(0,0700\text{ m}) $$
%             $$ v_i = 2,84 \cdot 10^6\text{ m/s} $$
            
%             \textbf{3. De tijd:}
%             Met EVRB-formule: $v_f = v_i + at$
%             $$ 0 = 2,84 \cdot 10^6\text{ m/s} + (-5,75 \cdot 10^{13}\text{ m/s}^2)t $$
%             $$ t = 4,94 \cdot 10^{-8}\text{ s} $$
%         \end{solution}
    
% \end{exercise}
% %———————————————————————————————————————

\end{document}